\section{An explicit twelve-variable polynomial}
\label{sec:primitive}

All twelve supplied unknowns in this section are strictly positive:
\[
\boxed{J,F,\alpha,z,h,i,T,s,w,\eta,\zeta,t.}
\]
Every other letter denotes a substituted polynomial, not another
unknown.  Set
\begin{align*}
q&=\beta J+1,& C&=q-F-\alpha,& X&=wq,&Y&=sq^3,\\
E&=XY,&k&=\eta+\zeta,&c&=kY+\eta,&a&=E+Y,\\
H&=4a+3,&\Delta&=a^2+H=(a+2)^2-1.
\end{align*}
Define the packed rank and the two retained outer factors by
\begin{align}
R&=((q-1)(q-F)+\alpha)(q^2-1)+(M_C+qM_F^*)J,\label{eq:rank}\\
N_k&=k-hE-R,\label{eq:index}\\
N_t&=(K+w)C+q-F-z(q-1).\label{eq:transport}
\end{align}
The root data are
\begin{align*}
\ell&=EkY,& A&=1+\ell(\ell+k),& I&=ic^2,\\
D&=1+\Delta I^2,&Q&=(\Delta I)^2,&d&=Q-1,\\
v&=cT,&Z&=c+RD,&b&=X+ac,\\
G&=\Delta c^2+1,& e&=Q(Dv^2+Z^2)-1,&B_*&=2QvZ.
\end{align*}
The letter $B_*$ is a root coefficient; it is distinct from the fixed
compiler base $B$.

The separating weights are also explicit polynomials:
\begin{align}
g&=c+X,&u&=\ell+k,&C_*&=v^2+D+Z,\nonumber\\
r&=(a+2)I,&C_2&=v+c+Rr,&S&=u+g,\label{eq:weights}\\
\lambda&=QC_*S+1,&\mu&=2\lambda C_2+S+1.\nonumber
\end{align}

\subsection{A finite norm recipe, without division or square roots}

Introduce \emph{formal computational symbols} $\xi,\theta$ in the
free rank-four algebra
\[
\mathcal A=\Z[J,F,\alpha,z,h,i,T,s,w,\eta,\zeta,t]
 [\xi,\theta]/(\xi^2-D,\theta^2-A).
\]
The fixed ports are coefficients in this display.  The symbols
$\xi,\theta$ are eliminated by the next finite operations and are
not Diophantine variables.
In this algebra set
\begin{align}
\mathcal Z&=Ht+b-H\mu\xi-H\theta,\nonumber\\
\mathcal K&=d(\mathcal Z^2+G)-H^2\lambda^2(e-B_*\xi),\nonumber\\
\mathcal E&=\mathcal K^2-4d^2G\mathcal Z^2.\label{eq:initialnorm}
\end{align}
For $U=U_0+U_1\theta$ write
$\Norm_\theta(U)=U_0^2-AU_1^2$, reducing $\xi^2=D$ throughout.
For $V=V_0+V_1\xi$ write
$\Norm_\xi(V)=V_0^2-DV_1^2$.  Define
\begin{equation}
\mathcal R=\Norm_\xi\bigl(\Norm_\theta(\mathcal E)\bigr).
\label{eq:fullnorm}
\end{equation}
This is an ordinary integer polynomial.  The two monic reduction rules
make its computation finite and unambiguous.  The delivered JSON
circuits expand every addition, subtraction, and multiplication in this
recipe, including all operations with fixed numerals.

The proposed single equation is
\begin{equation}
\boxed{\mathcal P_{12}:=N_kN_t(1+\mathcal R)-1=0.}
\label{eq:P12}
\end{equation}
Two useful equivalent presentations are
\begin{equation}
N_k=1,\qquad N_t=1,\qquad\mathcal R=0
\label{eq:system12}
\end{equation}
and the transparent sum of squares
\begin{equation}
(N_k-1)^2+(N_t-1)^2+\mathcal R^2=0.
\label{eq:sos12}
\end{equation}

\subsection{The preliminary bounds do not use the deleted roots}

Suppose the outer factors are integer units.  The identity
\[
q-F-z(q-1)=1-F-(z-1)(q-1)\le0
\]
and $K+w\ge2$ show that $C\le-1$ would force $N_t\le-2$.
Thus $C\ge0$ and $q-F=C+\alpha>0$.  The authentic mask bounds give
\begin{gather}
q\le(q-1)(q-F)+\alpha\le q^2-q,\nonumber\\
0<(M_C+qM_F^*)J<(q-1)(1+2q),\nonumber\\
q(q^2-1)<R<q^4-q^3+q^2-1<q^4-2.\label{eq:weakbounds}
\end{gather}
Since $X\ge q\ge16$ and $Y\ge q^3$, we have $E\ge q^4>R+2$.
All root data that must be positive are now positive, in particular
$c>1$, $I>1$, $D>1$, $Q>1$, and $d>0$.
This argument does not presuppose a main Pell root, a first Pell root,
or a solution of the strong equation.

\subsection{A rational auxiliary ordinate is automatically integral}

Write $\chi_u(n),\psi_u(n)$ for the integer Pell sequences specified by
\[
(u+\sqrt{u^2-1})^n=\chi_u(n)+\psi_u(n)\sqrt{u^2-1}.
\]

\begin{lemma}[First-coordinate divisibility]
\label{lem:firstdiv}
For $u>1$ and $m,n\ge1$,
$\chi_u(n)\mid\chi_u(m)$ implies $m=(2j+1)n$ for some $j\ge0$.
\end{lemma}
\begin{proof}
The addition formulas give
\[
\chi_u(2n+j)=2\chi_u(n)\chi_u(n+j)-\chi_u(j),
\]
and the corresponding reflection identity with $j$ replaced by $-j$.
Reduction of $m$ modulo $2n$, followed by reflection, makes
$\chi_u(m)$ congruent to $\pm\chi_u(r)$ modulo $\chi_u(n)$ for some
$0\le r\le n$.  Strict increase gives
$1\le\chi_u(r)<\chi_u(n)$ if $r<n$.  Divisibility forces $r=n$,
which is exactly the asserted odd-multiple condition.
\end{proof}

\begin{lemma}[Rational-ordinate rigidity]
\label{lem:rigidity}
Let $a_0>1$ and $V$ be integers.  If $y\in\Q$ satisfies
\begin{equation}
a_0^2V^2-(a_0^2-1)y^2=1,
\label{eq:rigidity}
\end{equation}
then $y\in\Z$.
\end{lemma}
\begin{proof}
Write $a_0^2-1=d_0s_0^2$ with $d_0>1$ squarefree and $s_0>0$.
For $r=s_0y$, the number $d_0r^2=(a_0V)^2-1$ is integral.
The square of the reduced denominator of $r$ divides squarefree $d_0$;
hence $r$ is integral.

Let $u+v_0\sqrt{d_0}$ be the fundamental positive solution of the
integer Pell equation.  Classification of its positive solutions gives
\begin{align*}
a_0&=\chi_u(n),&s_0&=v_0\psi_u(n),\\
|a_0V|&=\chi_u(m),&|r|&=v_0\psi_u(m)
\end{align*}
with $m,n\ge1$.  The zero-index cases are excluded by $a_0>1$ and
equation~\eqref{eq:rigidity}.  Lemma~\ref{lem:firstdiv} gives
$m=(2j+1)n$.  Pell composition then gives
\[
\psi_u(m)=\psi_u(n)\psi_{a_0}(2j+1),
\]
so $|r|=s_0\psi_{a_0}(2j+1)$ and $|y|$ is an integer.
\end{proof}

This special lemma is necessary: a rational root of a general nonmonic
integer quadratic need not be integral.  The main quotient $\gamma$
will be recovered by a different integrality argument.

\subsection{Separation of the sixteen branches}

Over an algebraic closure of $\Q$, consider the tuples
$(\gamma,\tau,y,f)$ satisfying
\begin{equation}
(H\gamma+b)^2=G,\qquad \tau^2=A,\qquad f^2=D,
\qquad dy^2=e-B_*f.
\label{eq:branches}
\end{equation}
There are exactly sixteen tuples.  The last radicand cannot vanish:
if $e-B_*f=0$, then $f=e/B_*\in\Q$ and $f^2=D$ makes $f$ integral.
But the definition of $e,B_*$ would give $Q(vf-Z)^2=1$, impossible
for integers with $Q>1$.  The same reasoning shows
\begin{equation}
e^2-B_*^2D\ne0.
\label{eq:nonzero-conjugatenorm}
\end{equation}
The auxiliary coordinate $y$ may be complex during this intermediate
argument.  Every bound below is an absolute-value bound and remains
valid for that case.

For every tuple~\eqref{eq:branches},
\begin{equation}
|\gamma|<g,\quad |\tau|<u,\quad \sqrt A>g,
\qquad \frac1{QC_*}<|y|,
\qquad \frac{|y|}{\sqrt D}<2C_2.
\label{eq:separationbounds}
\end{equation}
Here are the details.  Because $c>1$,
$\sqrt G<(a+2)c$, so
\[
|\gamma|<\frac{(2a+2)c+X}{4a+3}<c+X=g.
\]
Also $u^2-A=\ell k+k^2-1>0$, and
$\sqrt A>\ell=kXY^2>c+X$, using $c<k(Y+1)$,
$X\ge1$, $Y\ge2$, and $k\ge2$.

The arithmetic-geometric mean bound gives $v\sqrt D+Z<C_*$, hence
$e+B_*\sqrt D<QC_*^2$.  The two conjugates
$e\pm B_*\sqrt D$ have nonzero integer product by
\eqref{eq:nonzero-conjugatenorm}, and neither has absolute value greater
than $e+B_*\sqrt D$.  Each therefore has absolute value exceeding
$1/(QC_*^2)$.  Dividing by $d<Q$ proves the lower bound for $|y|$.
For the upper bound, $Q/d<2$ gives
\[
\frac{|y|}{\sqrt D}<\sqrt2\left(v+\frac Z{\sqrt D}\right).
\]
Finally $r^2-D=I^2-1>0$, so
$Z/\sqrt D<c+Rr$, proving the last inequality in
\eqref{eq:separationbounds}.

\begin{lemma}[Explicit primitive-element separation]
\label{lem:separate}
The sixteen values
\begin{equation}
\gamma+\tau+\lambda y+\mu f
\label{eq:compressed}
\end{equation}
are pairwise distinct.
\end{lemma}
\begin{proof}
Compare distinct tuples in the order $f,y,\tau,\gamma$.
If $f$ differs, its contribution has absolute value $2\mu\sqrt D$.
The sum of all lower-coordinate differences has absolute value less
than $4\lambda C_2\sqrt D+2S$, which is smaller by the definition of
$\mu$.  If $f$ agrees but $y$ differs, its weighted difference exceeds
$2\lambda/(QC_*)>2S$, whereas the remaining differences total less
than $2S$.  If only $\tau$ and possibly $\gamma$ differ, their respective
bounds are $2\sqrt A>2g$ and less than $2g$.  Finally the two main
roots give distinct $\gamma$ when every other coordinate agrees.
\end{proof}

\subsection{Norm factorization and integral root recovery}

Multiplication of the two main-root factors, then the two auxiliary-root
factors, gives the following exact factorization:
\begin{equation}
\mathcal R=d^8H^{16}
 \prod_{\text{sixteen tuples in \eqref{eq:branches}}}
 (t-\gamma-\tau-\lambda y-\mu f).
\label{eq:16factorization}
\end{equation}
For the scale, the two $\gamma$ factors multiply to
$((\mathcal Z-H\lambda y)^2-G)/H^2$.
Multiplication over the two signs of $y$, with $dy^2=e-B_*f$,
and multiplication by $d^2H^4$ yields exactly
$\mathcal E$ in~\eqref{eq:initialnorm}.  The two remaining quadratic
norms give~\eqref{eq:16factorization}.  In particular, $\mathcal R$
has degree sixteen in $t$, with leading coefficient $d^8H^{16}$.

Suppose $t\in\Z$ and $\mathcal R=0$.  Since $dH\ne0$, one tuple has
$t=\gamma+\tau+\lambda y+\mu f$.  Every automorphism of the algebraic
closure fixing $\Q$ permutes the sixteen tuples and fixes this sum.
Lemma~\ref{lem:separate} forces it to fix the entire tuple.  In
characteristic zero an algebraic number fixed by all these automorphisms
is rational.  Thus all four recovered coordinates are rational.

Since $f^2=D$ and $\tau^2=A$ are integers, $f,\tau\in\Z$.
Now $V=vf-Z$ is an integer, and the last equation of
\eqref{eq:branches} becomes
\[
QV^2-(Q-1)y^2=1.
\]
Use Lemma~\ref{lem:rigidity} with $a_0=\Delta I$ to obtain
$y\in\Z$.  Finally,
\[
\gamma=t-\tau-\lambda y-\mu f\in\Z.
\]
This last step retains the original congruence represented by the
main-root quotient; merely knowing $H\gamma+b\in\Z$ would not suffice.

The four deleted equations have now been restored, with possibly signed
roots.  Moreover $V=vf-Z=c(Tf-1)-Rf^2$ because $f^2=D$.
Setting $i_{\rm ord}=\Delta i$ gives
\[
(i_{\rm ord}c^2)^2=Q=\Delta(f^2-1).
\]
The signed ordinary compiler theorem described in
Section~\ref{sec:overview} therefore applies.  With $N_k=N_t=1$ it
proves soundness of~\eqref{eq:system12}.  Conversely, a canonical
positive normalized compiler witness gives a positive
$t=\gamma+\tau+\lambda y+\mu f$ and satisfies all three equations.

\subsection{The product finalizer is sound}

In the rank-four algebra modulo four, $\mathcal E\equiv\mathcal K^2$.
Quadratic norms are multiplicative over every commutative coefficient
ring, as follows directly from
\[
(ac+\delta bd)^2-\delta(ad+bc)^2
=(a^2-\delta b^2)(c^2-\delta d^2).
\]
Taking both norms shows that $\mathcal R$ is a square modulo four.
Consequently $1+\mathcal R$ cannot be $-1$.

If~\eqref{eq:P12} vanishes, its three integer factors are units.
Thus $\mathcal R=0$ and $N_kN_t=1$, although the two outer factors
may initially both be $-1$.  The weak bounds
\eqref{eq:weakbounds} permit both signs, so every preceding root
recovery remains valid.  The resulting tuple makes the full signed
ordinary core product equal one.  Its theorem proves blank halting
and recovers $N_k=N_t=1$ afterward.  No positive sign was assumed
before its justification.  Positive canonical witnesses prove the
converse, completing Theorem~\ref{thm:main}.

\subsection{Uniform exact degree and arithmetic size}

Assign degree one to each supplied unknown and degree zero to each
fixed port.  The relevant degrees are
\begin{center}\small
\begin{tabular}{c|rrrrrrrrrrr}
Quantity &$c$&$a,H$&$\Delta$&$R$&$\ell$&$D$&$Q$&$Z$&$e$&$\lambda$&$\mu$\\
\hline
Degree &5&6&12&4&11&34&46&38&122&95&116
\end{tabular}
\end{center}
Let a bar denote the highest homogeneous form.  Then
$\bar D=\bar a^2\bar I^2$, $\bar Z=\bar R\bar D$,
$\bar\lambda=\bar Q\bar R\bar D\bar\ell$, and
$\bar\mu=2\bar\lambda\bar R\bar a\bar I$.
In $\mathcal K$, the two degree-324 contributions have coefficient
ratio four to one, so they leave
\[
\overline{\mathcal K}
=3\bar Q\bar H^2\bar\lambda^2\bar R^2\bar D^2.
\]
All terms involving a surviving formal root have lower weighted
degree, with weights $\deg\xi=17$, $\deg\theta=11$.
The subtractand $4d^2G\mathcal Z^2$ has degree only 392, while
$\mathcal K^2$ has degree 648.  The two norms therefore give
\[
\overline{\mathcal R}
=3^8\bar Q^8\bar H^{16}\bar\lambda^{16}\bar R^{16}\bar D^{16}
=3^8\bar Q^{24}\bar H^{16}\bar R^{32}\bar D^{32}\bar\ell^{16}.
\]
This is nonzero for every valid coefficient tuple: in particular
$\bar R=(\beta J)^3(\beta J-F)$ is not the zero polynomial.
Thus $\deg\mathcal R=2592$, and
$\deg\mathcal P_{12}=7+2+2592=2601$.
The unique highest square in~\eqref{eq:sos12} gives degree 5184.
The generated circuits count respectively
\[
163=102M+61A,\qquad167=104M+63A,\qquad170=105M+65A,
\]
where $M$ counts multiplication and $A$ addition or subtraction.
